\section{\texorpdfstring{$K$}{K}-代数}\label{sec:k-algebras}

哪些类型的环可以写成某个仿射簇 $X$ 的 $K[X]$? 先考察基域 $K$ 所起的特殊作用, 从而给出这个问题的第一个答案.

为了引出下面的讨论, 考虑 $X=\mathbb{A}^{1}$ 的情形, 此时
\[
  K[\mathbb{A}^{1}]=K[x].
\]
如我们已经详细讨论的, $K[x]$ 是一个环. 不过, 它不只是一个普通的环, 还是 $K$ 上的向量空间, 因为除了能对多项式作加法与乘法, 也能用 $K$ 中的标量乘多项式 (加法和数乘这两个运算满足通常的向量空间公理). 与线性代数入门课程中通常研究的向量空间不同, $K[x]$ 是无限维的, 它的一组基为
\[
  B=\{1,x,x^{2},x^{3},\ldots\}.
\]
尽管如此, 与较熟悉的有限维向量空间一样, $K[x]$ 的每个元素都能唯一地写成 $B$ 中元素的线性组合.

更一般地, 每个坐标环都具有相同的附加结构: 它既是一个环, 又是 $K$ 上的向量空间. 本节建立 $K$-代数的代数基础, 用 $K$-代数将这一结构形式化. (在此提醒读者, 本书假设所有环都是带单位元的交换环.)

\begin{Definition}{$K$-代数}\leavevmode\label{def:k-algebra}
  一个 \textbf{$K$-代数}是一个环 $A$, 连同一个数乘映射
  \[
    \begin{aligned}
      K\times A&\to A\\
      (r,a)&\mapsto r\cdot a
    \end{aligned}
  \]
  满足以下公理.
  \begin{enumerate}[label=\arabic*.]
    \item 对所有 $r\in K$ 和所有 $a,b\in A$, 有 $r\cdot(a+b)=r\cdot a+r\cdot b$.
    \item 对所有 $r,s\in K$ 和所有 $a\in A$, 有 $(r+s)\cdot a=r\cdot a+s\cdot a$.
    \item 对所有 $r,s\in K$ 和所有 $a\in A$, 有 $(rs)\cdot a=r\cdot(s\cdot a)$.
    \item 对所有 $a\in A$, 有 $1\cdot a=a$.
    \item 对所有 $r\in K$ 和所有 $a,b\in A$, 有 $r\cdot(ab)=(r\cdot a)b=a(r\cdot b)$.
  \end{enumerate}
\end{Definition}

\BookPairedNoteLeft
  {为便于读懂 Definition~\ref{def:k-algebra} 中的公理, $K$ 和 $A$ 各自内部的乘积用元素直接并写来表示, 记号\textquotedbl{}$\cdot$\textquotedbl{}只用于数乘.}
  {前四条公理规定 $A$ 构成 $K$ 上的向量空间, 第五条公理说明数乘与 $A$ 中的乘法运算如何相容.}
\noindent 由上述公理可知, 对任意 $a\in A$, 用 $0\in K$ 数乘 $a$ 得到 $0\in A$, 用 $-1\in K$ 数乘 $a$ 得到加法逆元 $-a\in A$. 这也是一般向量空间都具有的性质 (Exercise~\ref{ex:zero-and-negative-one-scalar-action}).

\begin{Example}{多项式环}\leavevmode\label{ex:polynomial-rings-as-k-algebras}
  尤其从代数几何的角度来看, 多项式环 $K[x_{1},\ldots,x_{n}]$ 是 $K$-代数的典型例子. 的确, 除了可以对多项式作加法与乘法, 还可以用 $K$ 中的标量乘多项式, 而 Definition~\ref{def:k-algebra} 中的公理很容易验证.
\end{Example}

\begin{Example}{坐标环}\leavevmode\label{ex:coordinate-rings-as-k-algebras}
  设 $X$ 是一个仿射簇. 坐标环 $K[X]$ 构成一个 $K$-代数. 对任意 $F\in K[X]$ 和 $r\in K$, 定义 $r\cdot F\in K[X]$ 为下式给出的函数:
  \[
    (r\cdot F)(a_{1},\ldots,a_{n})=r\cdot F(a_{1},\ldots,a_{n}),
  \]
  其中右边的乘法是 $K$ 中通常的乘法. 为验证 $r\cdot F$ 确实是 $K[X]$ 的元素, 注意对某个 $f\in K[x_{1},\ldots,x_{n}]$, 有 $F=f|_{X}$, 因而 $r\cdot F=(r\cdot f)|_{X}$. 由于 $r\cdot f\in K[x_{1},\ldots,x_{n}]$, 可见 $r\cdot F\in K[X]$. Definition~\ref{def:k-algebra} 中的公理同样容易验证.
\end{Example}

\begin{Example}{$K$ 的扩张环}\leavevmode\label{ex:extension-rings-as-k-algebras}
  若 $A$ 是以 $K$ 为子环的任意环, 那么 $A$ 自然是一个 $K$-代数, 其中数乘就是 $A$ 中通常的环乘法. 事实上, 在我们的假设下 (环是带单位元的交换环), 每个非平凡的 $K$-代数都以这种方式得到. 更确切地说, 给定任意 $K$-代数 $A\ne\{0\}$, 都有一个典范包含映射 $K\to A$. 经由这个包含映射把 $K$ 看成 $A$ 的子环, 数乘就与 $A$ 中通常的乘法等同起来 (Exercise~\ref{ex:canonical-copy-of-k-in-algebra}). 特别地, 每个非平凡的 $K$-代数都以典范的方式包含一份 $K$.
\end{Example}

\begin{Example}{不是 $K$-代数的例子}\leavevmode\label{ex:integers-not-k-algebra}
  根据 Example~\ref{ex:extension-rings-as-k-algebras}, 任何不包含 $K$ 的非零环都不是 $K$-代数. 例如, 由于 $\mathbb{Z}$ 不包含任何域, 它不是任何域 $K$ 上的 $K$-代数.
\end{Example}

要完整建立 $K$-代数的理论, 还必须指明其间适当的态射. 既然 $K$-代数是带有数乘运算的环, $K$-代数同态就是保持数乘的环同态.

\begin{Definition}{$K$-代数同态}\leavevmode\label{def:k-algebra-homomorphism}
  设 $A$ 和 $B$ 为 $K$-代数. 一个 \textbf{$K$-代数同态} $\phi:A\to B$ 是满足
  \[
    \phi(r\cdot a)=r\cdot\phi(a)
  \]
  的环同态, 其中等式对所有 $r\in K$ 和 $a\in A$ 成立. 若 $\phi$ 有一个同样为 $K$-代数同态的逆, 则称 $\phi$ 是一个 \textbf{$K$-代数同构}, 并记 $A\cong B$.
\end{Definition}

同构似乎比双射要求更强: 不仅逆函数应当存在, 它还必须是 $K$-代数同态. 不过, 与群, 环和域的情形一样, 若 $\phi$ 是双射同态, 那么它的逆自动是同态 (Exercise~\ref{ex:bijective-k-algebra-homomorphism-is-isomorphism}).

\begin{Example}{从多项式环出发的 $K$-代数同态}\leavevmode\label{ex:k-algebra-evaluation-homomorphisms}
  考虑求值映射
  \[
    \begin{aligned}
      \phi:\mathbb{R}[x,y]&\to\mathbb{R}\\
      f(x,y)&\mapsto f(2,3).
    \end{aligned}
  \]
  稍加思考便可看出 $\phi$ 是一个 $\mathbb{R}$-代数同态. 此外, 知道 $\phi$ 是一个 $\mathbb{R}$-代数同态后, 还可看出它完全由 $x$ 和 $y$ 的像决定. 例如, 一旦知道 $\phi(x)=2$ 和 $\phi(y)=3$, 利用 $\phi$ 是保持数乘的环同态这一事实, 可得
  \[
    \phi(5x^{2}y+2y+xy)=5(2)^{2}(3)+2(3)+(2)(3)=72.
  \]
  更一般地, 对任意 $K$-代数 $A$ 和子集 $\{a_{1},\ldots,a_{n}\}\subseteq A$, 存在唯一的 $K$-代数同态
  \[
    \phi:K[x_{1},\ldots,x_{n}]\to A
  \]
  满足对所有 $i$ 都有 $\phi(x_{i})=a_{i}$ (Exercise~\ref{ex:polynomial-ring-map-determined-by-generators}). 这表明, 给出一个 $K$-代数同态 $K[x_{1},\ldots,x_{n}]\to A$ 等价于选定 $a_{1},\ldots,a_{n}\in A$.
\end{Example}

与群和环一样, $K$-代数也有第一同构定理, 它是证明两个 $K$-代数同构的基本工具. 要叙述这一定理, 首先必须定义 $K$-代数的商与子代数.

由于 $K$-代数 $A$ 是一个环, 已知对任何理想 $I\subseteq A$, 都可以构造商环 $A/I$. 下面的结果表明, 商环 $A/I$ 自然地从 $A$ 继承一个 $K$-代数结构.

\begin{Proposition}{商 $K$-代数}\leavevmode\label{prop:quotient-k-algebra}
  设 $A$ 为 $K$-代数. 若 $I\subseteq A$ 是一个理想, 则商环 $A/I$ 是一个 $K$-代数, 其中数乘定义为
  \begin{equation}\label{eq:quotient-algebra-scalar-multiplication}
    r\cdot[a]=[r\cdot a].
  \end{equation}
\end{Proposition}

\begin{Proof}
  我们必须验证 \eqref{eq:quotient-algebra-scalar-multiplication} 给出的数乘是良定义的. 关键在于理想自动对数乘封闭: 若 $a\in I$ 且 $r\in K$, 则
  \[
    r\cdot a=r\cdot(1a)=(r\cdot1)a.
  \]
  其中 $1$ 是 $A$ 的单位元. 由于 $r\cdot1\in A$ 且 $a\in I$, 理想的吸收性蕴含 $(r\cdot1)a\in I$. 因此 $r\cdot a\in I$, 这验证了 $I$ 对数乘封闭.

  由此出发, 要说明数乘是良定义的, 假设 $[a]=[b]$. 则 $a-b\in I$. 利用 $I$ 吸收数乘这一事实, 得 $r\cdot a-r\cdot b=r\cdot(a-b)\in I$, 这表明 $[r\cdot a]=[r\cdot b]$.

  由于 $I$ 是理想, 我们已经知道 $A/I$ 是环. 所以只剩下 Definition~\ref{def:k-algebra} 中的五条公理需要验证. 它们容易由 $A$ 中相应公理的成立得出.
\end{Proof}

现在给出第一同构定理所需的最后一个要素.

\begin{Definition}{子代数}\leavevmode\label{def:k-subalgebra}
  设 $A$ 为 $K$-代数. \textbf{子代数} $B\subseteq A$ 是满足对所有 $r\in K$ 和 $b\in B$ 都有 $r\cdot b\in B$ 的子环.
\end{Definition}

换言之, $K$-代数的一个子集是子代数, 当且仅当它既是子环又对数乘封闭, 因而本身构成一个 $K$-代数. 一个自然的例子是把 $K$-代数 $K[x]$ 作为 $K[x,y]$ 的子代数.

\begin{Theorem}{$K$-代数的第一同构定理}\leavevmode\label{thm:first-isomorphism-theorem-k-algebras}
  若 $\phi:A\to B$ 是一个 $K$-代数同态, 则
  \begin{enumerate}[label=(\roman*)]
    \item $\operatorname{im}(\phi)$ 是 $B$ 的子代数,
    \item $\ker(\phi)$ 是 $A$ 的理想, 并且
    \item 函数
    \[
      \begin{aligned}
        [\phi]:\frac{A}{\ker(\phi)}&\to\operatorname{im}(\phi)\\
        [a]&\mapsto\phi(a)
      \end{aligned}
    \]
    是一个良定义的 $K$-代数同构.
  \end{enumerate}
\end{Theorem}

\begin{Proof}
  见 Exercise~\ref{ex:prove-first-isomorphism-theorem-k-algebras}.
\end{Proof}

\begin{Example}{$\mathbb{C}$ 作为商 $\mathbb{R}$-代数}\leavevmode\label{ex:complex-numbers-as-quotient-real-algebra}
  把 $\mathbb{C}$ 看成 $\mathbb{R}$-代数, 其中数乘是通常的乘法. 令 $\phi$ 为函数
  \[
    \begin{aligned}
      \phi:\mathbb{R}[x]&\to\mathbb{C}\\
      \phi(f)&=f(\mathrm{i}).
    \end{aligned}
  \]
  可以验证 (Exercise~\ref{ex:verify-complex-quotient-real-algebra}) $\phi$ 是以 $\langle x^{2}+1\rangle$ 为核的满 $\mathbb{R}$-代数同态. 因此, 由 Theorem~\ref{thm:first-isomorphism-theorem-k-algebras}, 有 $\mathbb{R}$-代数同构
  \[
    \mathbb{C}\cong\mathbb{R}[x]/\langle x^{2}+1\rangle.
  \]
  作为 $\mathbb{R}$ 上的向量空间, 复数集 $\mathbb{C}$ 有一组基 $\{1,\mathrm{i}\}$, 而商 $\mathbb{R}[x]/\langle x^{2}+1\rangle$ 有一组基 $\{1,[x]\}$. 同构 $[\phi]$ 将 $1$ 与 $1$ 等同, 将 $[x]$ 与 $\mathrm{i}$ 等同. 这一对应的由来是 $[x]^{2}=\mathrm{i}^{2}=-1\in\mathbb{R}$.
\end{Example}

回到引出这一讨论的坐标环, 现在可以用有关 $K$-代数的知识, 把 Proposition~\ref{prop:coordinate-ring-as-quotient} 改述为 $K$-代数情形下的结果. 设 $X\subseteq\mathbb{A}^{n}$ 为仿射簇, 考虑环同态
\[
  \begin{aligned}
    K[x_{1},\ldots,x_{n}]&\to K[X]\\
    f&\mapsto f|_{X}.
  \end{aligned}
\]
由于 $(r\cdot f)|_{X}=r\cdot(f|_{X})$, 这是一个 $K$-代数同态. 与 Proposition~\ref{prop:coordinate-ring-as-quotient} 的证明完全一样, 它是以 $\mathcal{I}(X)$ 为核的满射. 因此由第一同构定理得到如下结果.

\begin{Proposition}{坐标环的商 $K$-代数表示}\leavevmode\label{prop:coordinate-ring-as-quotient-k-algebra}
  若 $X\subseteq\mathbb{A}^{n}$ 是仿射簇, 则有典范 $K$-代数同构
  \[
    K[X]=\frac{K[x_{1},\ldots,x_{n}]}{\mathcal{I}(X)}.
  \]
\end{Proposition}

\BookExerciseHeading
\begin{BookExercises}
  \item\label{ex:zero-and-negative-one-scalar-action} 设 $A$ 是 $K$-代数. 证明对所有 $a\in A$, 都有 $0\cdot a=0$ 和 $(-1)\cdot a=-a$.
  \item\label{ex:bijective-k-algebra-homomorphism-is-isomorphism} 设 $\phi:A\to B$ 是 $K$-代数同态. 假设 $\phi$ 是双射, 因而有逆函数 $\phi^{-1}:B\to A$. 证明这个逆函数是 $K$-代数同态.
  \item\label{ex:polynomial-ring-map-determined-by-generators} 设 $A$ 是 $K$-代数, $\{a_{1},\ldots,a_{n}\}\subseteq A$ 为一个子集. 证明存在唯一的 $K$-代数同态 $\phi:K[x_{1},\ldots,x_{n}]\to A$, 使得对所有 $i=1,\ldots,n$ 都有 $\phi(x_{i})=a_{i}$.
  \item\label{ex:ring-homomorphism-not-complex-algebra-homomorphism} 举出一个不是 $\mathbb{C}$-代数同态的环同态 $\phi:\mathbb{C}\to\mathbb{C}$.
  \item\label{ex:canonical-copy-of-k-in-algebra} 设 $A\ne\{0\}$ 是一个 $K$-代数.
    \begin{enumerate}[label=(\alph*),ref=(\alph*)]
      \item\label{item:canonical-field-embedding} 证明由 $\phi(r)=r\cdot1$ 定义的函数 $\phi:K\to A$ 是单的环同态.
      \item 通过 \ref{item:canonical-field-embedding} 中的同态把 $K$ 看作 $A$ 的子环, 证明数乘与 $A$ 中通常的乘法等同.
    \end{enumerate}
  \item\label{ex:k-algebra-homomorphism-restricts-to-identity} 设 $A$ 和 $B$ 是非平凡的 $K$-代数. 由 Exercise~\ref{ex:canonical-copy-of-k-in-algebra}, 两者都以典范的方式包含一份 $K$. 证明任何 $K$-代数同态 $\phi:A\to B$ 限制到 $K$ 上都是恒等映射.
  \item\label{ex:prove-first-isomorphism-theorem-k-algebras} 证明 $K$-代数的第一同构定理.
  \item\label{ex:verify-complex-quotient-real-algebra} 补全 Example~\ref{ex:complex-numbers-as-quotient-real-algebra} 中的证明细节: 作为 $\mathbb{R}$-代数, $\mathbb{C}\cong\mathbb{R}[x]/\langle x^{2}+1\rangle$.
\end{BookExercises}
